% ============================================================
\chapter{Giới thiệu}
\label{chap:intro}
% ============================================================

\section{Bối cảnh và động lực nghiên cứu}
\label{sec:intro:motivation}

Sự bùng nổ của Internet vạn vật (IoT) và những bước tiến vượt bậc trong triển khai mạng
không dây thế hệ thứ năm (5G) và thứ sáu (6G) đang đặt ra yêu cầu ngày càng khắt khe về
khả năng kết nối đáng tin cậy, độ trễ thấp và hiệu quả năng lượng cho hàng tỷ thiết bị
đầu cuối. Trong bối cảnh đó, mô hình truyền thông thiết bị-với-thiết bị (Device-to-Device
--- D2D) nổi lên như một giải pháp kiến trúc then chốt: thay vì định tuyến mọi lưu lượng
qua trạm gốc, các thiết bị liền kề có thể giao tiếp trực tiếp với nhau, từ đó giảm tải
mạng lõi, tối ưu hiệu suất phổ tần cục bộ và cải thiện trải nghiệm người dùng
\cite{gu2023uav}.

Tuy nhiên, môi trường vô tuyến vốn mang đặc tính mở và chia sẻ, tạo điều kiện thuận lợi
cho các cuộc tấn công gây nhiễu (jamming) có chủ đích. Trong kiểu tấn công này, kẻ địch
phát liên tục một tín hiệu RF công suất cao nhằm áp đảo tín hiệu hợp lệ, làm suy giảm
nghiêm trọng tỷ số tín hiệu trên nhiễu và tạp âm (SINR) tại các thiết bị thu, dẫn đến
sụt giảm thông lượng, tăng tỷ lệ lỗi bit, và trong trường hợp cực đoan có thể vô hiệu hóa
hoàn toàn đường truyền \cite{pelechrinis2011dos}. Các thiết bị IoT -- vốn bị ràng buộc
chặt chẽ về dung lượng pin, khả năng xử lý và không gian lưu trữ -- đặc biệt dễ bị tổn
thương trước kiểu tấn công này.

Các phương pháp phòng thủ truyền thống như nhảy tần (Frequency Hopping) hay trải phổ
(Spread Spectrum) đang ngày càng bộc lộ giới hạn. Khi đối mặt với các jammer thông minh
có khả năng theo dõi và bắt chước chiến lược của hệ thống hợp lệ, những phương pháp phản
ứng tĩnh này không còn đủ hiệu quả \cite{grover2014jamming}. Quan trọng hơn, tất cả các
tiếp cận này đều coi tín hiệu jamming là mối đe dọa cần né tránh, trong khi bản thân tín
hiệu RF công suất cao từ jammer lại chứa đựng một nguồn năng lượng miễn phí đáng kể.

Nhận thức đó thúc đẩy việc ứng dụng kỹ thuật \emph{Ambient Backscatter} (AmBC) vào bài
toán chống jamming. Thay vì phát sóng tự thân -- đòi hỏi nguồn điện riêng -- các thiết bị
backscatter chỉ cần điều chế và phản xạ lại tín hiệu RF môi trường (ở đây chính là tín
hiệu jamming) để truyền thông tin đến đích. Công suất tiêu thụ của thẻ backscatter gần như
bằng không, biến AmBC trở thành ứng viên lý tưởng cho các thiết bị IoT giới hạn năng lượng
\cite{huynh2019jam}. Điểm cốt lõi trong đề tài này là một sự đảo ngược quan niệm: hệ thống
không coi jamming là kẻ thù, mà chủ động \emph{biến kẻ thù thành tài nguyên}, khai thác
chính công suất RF của jammer làm nguồn kích hoạt cho truyền thông.

Mặc dù vậy, do tín hiệu backscatter có mức công suất rất thấp, chất lượng đường truyền D2D
phụ thuộc mạnh vào khoảng cách và điều kiện kênh truyền. Khi hai thiết bị ở xa nhau hoặc bị
che khuất bởi địa hình -- tình huống phổ biến trong các môi trường đô thị và công nghiệp --
tín hiệu phản xạ không đủ mạnh để truyền tin cậy. Để giải quyết nút thắt này, đề tài khai
thác tính linh hoạt ba chiều của thiết bị bay không người lái (UAV) với vai trò relay trung
gian động. Nhờ khả năng di chuyển tự do trong không gian 3D, UAV có thể chủ động bay đến
vị trí tối ưu để đồng thời cải thiện cả hai chặng kênh truyền (nguồn-to-UAV và UAV-to-đích),
đồng thời tận dụng lợi thế xác suất đường truyền thẳng (LoS) cao của kênh không-đất (A2G)
\cite{duo2020antijamming3d}.

Điều phối nhiều UAV và nhiều thiết bị SU để đưa ra quyết định tối ưu trong thời gian thực
là bài toán đa tác tử có độ phức tạp cao: không gian trạng thái liên tục, thông tin quan sát
không đầy đủ (do mỗi tác tử chỉ có thể quan sát trạng thái cục bộ), và môi trường phi tuyến
thay đổi liên tục. Học tăng cường sâu đa tác tử (MARL), đặc biệt là MADDPG với paradigm
huấn luyện tập trung -- thực thi phân tán (CTDE), đã chứng minh hiệu quả trong các bài toán
tương tự \cite{lowe2017maddpg}. Tuy nhiên, MADDPG tiêu chuẩn bị hạn chế bởi bất ổn
cold-start và khó hội tụ trong không gian hành động nhiều chiều, đòi hỏi những kỹ thuật
bổ sung để đảm bảo chất lượng huấn luyện.

Từ những phân tích trên, đề tài đặt ra câu hỏi nghiên cứu trung tâm:

\begin{quote}
\itshape
Làm thế nào để huấn luyện đồng thời nhiều UAV agent và nhiều SU agent -- sử dụng học tăng
cường sâu kết hợp imitation learning và Transformer-GAT -- nhằm tối đa hóa thông lượng D2D,
đáp ứng ràng buộc QoS và tiết kiệm năng lượng, trong điều kiện bị tấn công gây nhiễu liên
tục và thông tin kênh truyền không đầy đủ?
\end{quote}

\section{Tình hình nghiên cứu trong và ngoài nước}
\label{sec:intro:sota}

\subsection{Nghiên cứu ngoài nước}
\label{subsec:intro:sota:global}

\subsubsection*{Gây nhiễu và chống gây nhiễu trong mạng không dây}

Grover và cộng sự~\cite{grover2014jamming} trình bày khảo sát toàn diện về các kỹ thuật
jamming và anti-jamming trong mạng không dây, bao gồm phân loại hình thức tấn công (spot,
sweep, barrage, deceptive jammer), đánh giá tác động lên hiệu năng hệ thống và so sánh các
giải pháp phòng thủ từ lớp vật lý đến giao thức MAC. Pelechrinis và cộng sự~\cite{pelechrinis2011dos}
phân tích sâu bản chất của các tấn công DoS do jamming, cung cấp tiêu chí phân loại và khung
so sánh hiệu quả các kỹ thuật phòng thủ, đặt nền tảng học thuật cho nhiều nghiên cứu tiếp theo.

\subsubsection*{Ứng dụng học tăng cường sâu trong chống gây nhiễu}

Van Huynh và cộng sự~\cite{huynh2019jam} là nhóm tiên phong đề xuất kết hợp Ambient
Backscatter với Deep Dueling Neural Network để chống jamming. Phương pháp cho phép hệ thống
học chiến lược lựa chọn kênh và điều chỉnh hệ số phản xạ tối ưu thông qua tương tác với
môi trường, đạt được hiệu năng vượt trội so với các phương pháp truyền thống. Liu và cộng
sự~\cite{liu2018antijamming} ứng dụng DRL với đặc trưng hình ảnh phổ thay đổi theo thời
gian-tần số (spectrum waterfall) để nhận diện và thích ứng với các dạng jamming đa dạng trong
thời gian thực. Dang và cộng sự~\cite{dang2024ppo} đề xuất thuật toán PPO (Proximal Policy
Optimization) giúp hệ thống tự học chính sách lựa chọn kênh và chiến lược điều chế tối ưu
trong môi trường jamming biến đổi liên tục. Mặc dù đạt kết quả tốt, ba nghiên cứu này đều
giải quyết bài toán đơn tác tử (single-agent), chưa khai thác được khả năng hợp tác phối
hợp giữa nhiều thiết bị.

\subsubsection*{UAV trong mạng truyền thông không dây}

Gu và Zhang~\cite{gu2023uav} cung cấp khảo sát toàn diện về ứng dụng UAV hỗ trợ truyền
thông không dây, bao gồm tối ưu quỹ đạo, phân bổ tài nguyên, thiết kế giao thức và các xu
hướng nghiên cứu mới nhất. Wang và cộng sự~\cite{wang2023uavjamming} phân tích khả năng
chống jamming của hệ thống UAV, đề xuất mô hình vùng phần thưởng tương tác (IRR) cho phép
tối ưu đồng thời công suất phát và quỹ đạo bay, mở rộng phạm vi bảo vệ hợp lệ. Ma và cộng
sự~\cite{ma2022rl} đề xuất khung điều khiển công suất chống jamming động dựa trên DDPG, ước
lượng cường độ tín hiệu nhiễu hiệu quả bằng Kernel Density Estimation (KDE), đồng thời thiết
kế quỹ đạo UAV theo phân cụm người dùng K-means++. Duo và cộng sự~\cite{duo2020antijamming3d}
tập trung thiết kế quỹ đạo 3D chống jamming cho mạng cảm biến không dây với kênh LoS xác
suất, tối đa hóa tốc độ kỳ vọng tối thiểu qua thuật toán tối ưu lồi liên tiếp. Các phương
pháp trên tuy đạt kết quả tốt nhưng đều dựa vào đơn UAV đơn tác tử hoặc quy tắc tĩnh, thiếu
khả năng học hợp tác đa UAV.

Liên quan trực tiếp nhất đến đề tài là công trình IA-MADDPG gốc~\cite{iamaddpg2024}, đề xuất
framework kết hợp imitation learning với MADDPG cho mạng D2D AmBC chống jamming với relay
RBS cố định. Công trình chứng minh hiệu quả rõ rệt của việc warm-up replay buffer bằng expert
policy và áp dụng behavior cloning để ổn định quá trình hội tụ trong bài toán đa tác tử, mở
ra nền tảng lý thuyết và thực nghiệm cho hướng mở rộng mà đề tài này theo đuổi.

\subsection{Nghiên cứu trong nước}
\label{subsec:intro:sota:vietnam}

Tại Việt Nam, nghiên cứu về mạng UAV hiện tập trung chủ yếu vào các ứng dụng thực tiễn:
giám sát biên giới và quản lý không phận, hỗ trợ tìm kiếm cứu nạn, giám sát hạ tầng và
nông nghiệp thông minh \cite{hang2023uav, vohoa2023uav}. Viettel High Tech đã triển khai
nhiều mẫu UAV trinh sát phục vụ quốc phòng, trong khi các tổ chức như VNPT và Học viện Kỹ
thuật Quân sự nghiên cứu kết nối UAV với vệ tinh nhằm mở rộng phạm vi liên lạc trong địa
hình phức tạp \cite{viettel2025uav}.

Tuy nhiên, số lượng công bố quốc tế về tối ưu hóa kiến trúc mạng UAV, phân bổ tài nguyên
vô tuyến và bảo mật lớp vật lý trong bối cảnh Việt Nam còn rất hạn chế. Đặc biệt, chưa có
nghiên cứu trong nước nào đề cập đến việc kết hợp Ambient Backscatter, UAV relay và học
tăng cường đa tác tử trong một framework thống nhất chống gây nhiễu. Đây là khoảng trắng
nghiên cứu đáng chú ý, nhất là trong bối cảnh Việt Nam đang đẩy mạnh hạ tầng số quốc gia
và phát triển các công nghệ IoT nội địa.

\section{Khoảng trống nghiên cứu và đề xuất giải pháp}
\label{sec:intro:gap}

Phân tích tổng quan tài liệu làm rõ bốn khoảng trống nghiên cứu chính:

\textbf{Khoảng trống thứ nhất --- Mô hình đơn tác tử.}
Hầu hết các nghiên cứu DRL chống jamming cho UAV \cite{wang2023uavjamming, ma2022rl,
duo2020antijamming3d} chỉ xét bài toán một UAV (single-agent), bỏ qua khả năng hợp tác và
phân công vai trò giữa nhiều UAV và nhiều thiết bị SU trong mạng lớn hơn. Khi số lượng cặp
SU--DU tăng lên, một UAV duy nhất không thể phục vụ hiệu quả toàn bộ mạng.

\textbf{Khoảng trống thứ hai --- Chưa khai thác năng lượng jamming.}
Các nghiên cứu UAV chống jamming hiện tại \cite{wang2023uavjamming, ma2022rl,
duo2020antijamming3d} coi tín hiệu jamming thuần túy là nhiễu cần né tránh, chưa tích hợp
cơ chế Ambient Backscatter để biến năng lượng RF của jammer thành tài nguyên truyền thông.
Sự kết hợp UAV relay với AmBC tạo ra một kiến trúc hệ thống hoàn toàn mới chưa được khám
phá đầy đủ trong tài liệu.

\textbf{Khoảng trống thứ ba --- Bất ổn cold-start trong MARL.}
MADDPG tiêu chuẩn bị bất ổn cold-start nghiêm trọng khi huấn luyện trong không gian hành
động liên tục nhiều chiều với nhiều tác tử dị thể (SU và UAV có không gian hành động và
quan sát khác nhau). Imitation learning kết hợp với MARL dị thể vẫn là hướng nghiên cứu
còn nhiều tiềm năng khám phá.

\textbf{Khoảng trống thứ tư --- Critic không khai thác cấu trúc đồ thị.}
Critic MLP truyền thống trong MADDPG coi toàn bộ quan sát như một vector phẳng, không nắm
bắt được cấu trúc topo đồ thị kết nối UAV--SU (biến đổi theo thời gian khi UAV di chuyển)
và tương quan không gian giữa các tác tử. Điều này dẫn đến ước lượng giá trị Q kém chính
xác, làm chậm tốc độ hội tụ.

Để lấp đầy các khoảng trống trên, đề tài đề xuất framework \textbf{IA-MADDPG với
Transformer-GAT Centralized Critic} cho bài toán mạng D2D AmBC chống jamming có hỗ trợ
UAV relay động. Framework tích hợp đồng thời: (i) $K = 2$ UAV relay di động 3D bên cạnh
RBS cố định; (ii) khai thác năng lượng jamming qua AmBC với ba chế độ truyền linh hoạt;
(iii) expert policy để warm-up replay buffer và behavior cloning để ổn định giai đoạn đầu
huấn luyện; (iv) Transformer encoder xử lý quan sát đa tác tử dưới điều kiện thông tin
không đầy đủ; và (v) Graph Attention Network mã hóa cấu trúc đồ thị UAV--SU trong
centralized critic.

\section{Đóng góp của đề tài}
\label{sec:intro:contributions}

Đề tài đóng góp năm điểm mới có giá trị học thuật và thực tiễn:

\begin{enumerate}

  \item \textbf{Mô hình hệ thống tích hợp UAV relay động và Ambient Backscatter.}
  Đề xuất và phân tích mô hình mạng D2D AmBC gồm $N$ cặp SU--DU mặt đất, $K$ UAV relay
  di động, một RBS cố định và một jammer mặt đất liên tục trong khu vực $200 \times 200\,$m$^2$.
  Hệ thống khai thác tín hiệu jamming làm RF excitation cho Ambient Backscatter, với ba chế
  độ truyền linh hoạt cùng mô hình kênh Rayleigh mặt đất và kênh LoS xác suất không-đất.

  \item \textbf{Phát biểu bài toán Dec-POMDP đa tác tử dị thể.}
  Công thức hóa bài toán tối ưu liên hợp --- hệ số phản xạ $\alpha_i$, lựa chọn chế độ
  truyền $m_i$ (rời rạc) và quỹ đạo UAV 3D $\Delta\mathbf{p}_{U_k}$ (liên tục) --- dưới
  dạng Dec-POMDP với $N + K$ tác tử dị thể, không gian quan sát cục bộ khác nhau và không
  gian hành động lai (liên tục và rời rạc hỗn hợp).

  \item \textbf{Transformer-GAT Centralized Critic.}
  Thiết kế kiến trúc critic mới thay thế MLP truyền thống: bộ mã hóa Transformer khai thác
  cơ chế self-attention để tổng hợp ngữ cảnh toàn cục từ quan sát của tất cả tác tử trong
  điều kiện thông tin không đầy đủ; Graph Attention Network (GAT) mã hóa cấu trúc đồ thị
  kết nối UAV--SU động, nắm bắt các phụ thuộc không gian và topo mạng tức thời.

  \item \textbf{Expert policy cho UAV agent và giao thức warm-up IA-MADDPG mở rộng.}
  Xây dựng thuật toán heuristic vị trí UAV dựa trên thông tin kênh đầy đủ (full CSI), làm
  expert policy cho giai đoạn warm-up replay buffer, kết hợp với behavior cloning loss có
  hệ số annealing $\lambda_{\mathrm{IL}}$ để ổn định huấn luyện và đạt hội tụ nhanh hơn so
  với MADDPG tiêu chuẩn.

  \item \textbf{Đánh giá toàn diện và nghiên cứu ablation.}
  Thực hiện mô phỏng Python so sánh định lượng với năm phương pháp baseline (truyền thẳng
  DT, Greedy SINR, nhảy tần FH, MADDPG tiêu chuẩn, IA-MADDPG gốc không có UAV) trên các
  chỉ số phần thưởng mạng, TSR, thông lượng, hiệu quả năng lượng và quỹ đạo UAV; đồng thời
  thực hiện nghiên cứu ablation định lượng đóng góp độc lập của từng thành phần đề xuất.

\end{enumerate}

\section{Cấu trúc báo cáo}
\label{sec:intro:structure}

Phần còn lại của báo cáo được tổ chức thành năm chương như sau.

\textbf{Chương~\ref{chap:background} --- Cơ sở lý thuyết} trình bày các kiến thức nền tảng
cần thiết, bao gồm nguyên lý và phân loại truyền thông Backscatter (monostatic, bistatic,
ambient), mô hình tấn công jamming và các kỹ thuật chống nhiễu truyền thống, mô hình kênh
không-đất của UAV với xác suất LoS, cùng các thuật toán học tăng cường từ DQN, DDPG, TD3
đến MADDPG, imitation learning, kiến trúc Transformer và Graph Attention Network.

\textbf{Chương~\ref{chap:system} --- Mô hình hệ thống và bài toán tối ưu} xây dựng chi tiết
mô hình mạng với tất cả thành phần, mô hình kênh Rayleigh và A2G xác suất LoS, biểu thức
SINR cho ba chế độ truyền, thiết kế hàm phần thưởng composite, và phát biểu chính thức bài
toán tối ưu dưới dạng Dec-POMDP với đầy đủ không gian quan sát, hành động và ràng buộc.

\textbf{Chương~\ref{chap:algorithm} --- Thuật toán IA-MADDPG với Transformer-GAT Critic}
trình bày kiến trúc chi tiết của framework đề xuất: cấu trúc Actor cho SU agent và UAV
agent, thiết kế Transformer-GAT Centralized Critic, quy trình hai giai đoạn (expert warm-up
và behavior cloning với annealing), các kỹ thuật hỗ trợ (PER, TD3 smoothing, xử lý không
gian hành động lai), pseudocode thuật toán đầy đủ và phân tích độ phức tạp.

\textbf{Chương~\ref{chap:simulation} --- Thiết lập mô phỏng và đánh giá hiệu năng} mô tả
môi trường mô phỏng Python, tham số hệ thống và huấn luyện, cách cài đặt năm phương pháp
baseline, và trình bày toàn bộ kết quả thực nghiệm bao gồm đường cong hội tụ, so sánh phần
thưởng, thông lượng, TSR, phân phối chế độ truyền, quỹ đạo UAV và nghiên cứu ablation.

\textbf{Chương~\ref{chap:conclusion} --- Kết luận và hướng phát triển} tổng kết các đóng
góp, phân tích hạn chế hiện tại và đề xuất các hướng nghiên cứu tiếp theo gồm: CSI không
hoàn hảo, jammer di động, đa jammer, nguồn năng lượng dị thể và triển khai quy mô lớn.
